% Source: https://artofproblemsolving.com/community/c3h1992167_a_formula_
% Stable identity: aops-topic-1992167-q01
(a) If $d$ divides $n$, then every prime factor of $d$ is one of the $p_i$, and $p_i$ cannot appear in $d$ more often than in $n$. So, by unique factorization, $d=p_1^{f_1}p_2^{f_2}\cdots p_k^{f_k}$ with $0\le f_i\le e_i$. Conversely every such product divides $n$, and different exponent choices give different numbers. Hence the divisors correspond exactly to the choices of $(f_1,\dots,f_k)$. There are $e_i+1$ choices for each $f_i$ (namely $0,1,\dots,e_i$), so by the multiplication principle there are $(e_1+1)(e_2+1)\cdots(e_k+1)$ divisors. $\blacksquare$

(b) $720=2^4\cdot3^2\cdot5$, so it has $(4+1)(2+1)(1+1)=\boxed{30}$ positive divisors.
